% =====================================================================
%  CARNET D'ENTRAÎNEMENT — Fiche 12 (Chimie organique)
%  THÈME 4 — Masse molaire d'une molécule organique
%  NIVEAU FACILE — ÉNONCÉS
% =====================================================================
\documentclass[11pt,a4paper]{article}
\usepackage[utf8]{inputenc}
\usepackage[T1]{fontenc}
\usepackage{lmodern}
\usepackage[french]{babel}
\usepackage{amsmath,amssymb,mathtools}
\providecommand{\degree}{^\circ}
\usepackage{icomma}
\usepackage{siunitx}
\sisetup{output-decimal-marker={,},per-mode=symbol}
\usepackage[version=4]{mhchem}
\usepackage{chemfig}
\setchemfig{atom sep=2em,bond length=0.9em,cram width=2.5pt}
\usepackage{xcolor}
\usepackage{array}
\usepackage{booktabs}
\usepackage{tabularx}
\usepackage[most]{tcolorbox}
\usepackage{enumitem}
\usepackage{tikz}
\usetikzlibrary{arrows.meta,positioning,calc}
\usepackage{fancyhdr}
\usepackage{titlesec}
\usepackage[hidelinks]{hyperref}
\usepackage[a4paper,margin=2.1cm,headheight=26pt,headsep=9pt]{geometry}

\definecolor{primary}{RGB}{150,98,162}
\definecolor{primarydark}{RGB}{92,52,110}
\definecolor{excol}{RGB}{0,135,90}
\definecolor{attncol}{RGB}{200,40,55}
\definecolor{methcol}{RGB}{90,90,100}

\newcommand{\runtitle}{Thème 4 — Facile — Énoncés}
\pagestyle{fancy}\fancyhf{}
\fancyhead[L]{\small\textcolor{primarydark}{\textbf{Fiche 12} — Chimie organique}}
\fancyhead[R]{\small\textcolor{primarydark}{\runtitle}}
\fancyfoot[C]{\small\thepage}
\renewcommand{\headrulewidth}{0.8pt}
\renewcommand{\headrule}{\hbox to\headwidth{\color{primary}\leaders\hrule height \headrulewidth\hfill}}
\setlength{\parindent}{0pt}\setlength{\parskip}{3pt}

\tcbset{boxbase/.style={enhanced,breakable,boxrule=0.9pt,arc=2.5pt,
   left=3mm,right=3mm,top=2mm,bottom=2mm,fonttitle=\bfseries,coltitle=white,
   toptitle=0.5mm,bottomtitle=0.5mm}}
\newtcolorbox[auto counter]{exo}[1]{boxbase,colback=primary!4,colframe=primary,
   title={\textsc{Exercice}~\thetcbcounter\ \textmd{--- \itshape #1}}}
\newtcolorbox{consigne}{boxbase,colback=methcol!8,colframe=methcol,sharp corners,
   title={\textsc{Consignes \& données}}}

\newcommand{\banner}[2]{%
\begin{tcolorbox}[enhanced,colback=primary,colframe=primary,arc=5pt,boxrule=0pt,
   left=5mm,right=5mm,top=2.5mm,bottom=2.5mm,coltext=white]
{\large\bfseries #1}\\[1pt]{\small #2}
\end{tcolorbox}\vspace{2pt}}

\begin{document}

\banner{Thème 4 — Masse molaire d'une molécule organique}{Niveau \textbf{Facile} — Énoncés \quad$\vert$\quad Fiche 12, Chimie organique}

\begin{consigne}
Exercices \textbf{ouverts}. Détaille toujours le calcul \emph{contribution par contribution}. \textbf{Données} : $M(\ce{H})=\qty{1}{\g\per\mol}$, $M(\ce{C})=\qty{12}{\g\per\mol}$, $M(\ce{N})=\qty{14}{\g\per\mol}$, $M(\ce{O})=\qty{16}{\g\per\mol}$.
\end{consigne}

% ---------------------------------------------------------------------
\begin{exo}{le calcul modèle}
Calcule la masse molaire $M$ du composé de formule brute $\ce{C6H14N2}$. Présente les trois contributions (C, H, N) puis leur somme.
\end{exo}

% ---------------------------------------------------------------------
\begin{exo}{deux petites molécules}
Calcule la masse molaire de :
\begin{enumerate}[label=\textbf{\alph*)},leftmargin=2em,itemsep=1pt]
  \item le méthane $\ce{CH4}$ ;
  \item l'éthanol $\ce{C2H6O}$.
\end{enumerate}
\end{exo}

% ---------------------------------------------------------------------
\begin{exo}{masse molaire ou masse moléculaire relative ?}
On a trouvé, pour une molécule, le nombre $114$.
\begin{enumerate}[label=\textbf{\alph*)},leftmargin=2em,itemsep=1pt]
  \item Comment écrit-on la \textbf{masse molaire} $M$ (avec son unité) ?
  \item Comment écrit-on la \textbf{masse moléculaire relative} $M_r$ ?
  \item Pourquoi la réponse «\,$M_r = \qty{114}{\g\per\mol}$\,» serait-elle fausse ?
\end{enumerate}
\end{exo}

% ---------------------------------------------------------------------
\begin{exo}{de la structure à la masse molaire}
Soit la propanone (acétone) :
\begin{center}\chemfig{CH_3-[:30](=[:90]O)-[:-30]CH_3}\end{center}
\begin{enumerate}[label=\textbf{\alph*)},leftmargin=2em,itemsep=1pt]
  \item Détermine sa formule brute (compte les C, H, O).
  \item Calcule sa masse molaire.
\end{enumerate}
\end{exo}

% ---------------------------------------------------------------------
\begin{exo}{de la masse à la quantité de matière}
On dispose de $\qty{9,2}{\g}$ d'éthanol ($M=\qty{46}{\g\per\mol}$). Calcule la \textbf{quantité de matière} $n$ correspondante (formule $n=m/M$).
\end{exo}

\end{document}
